\chapter{Wielomian Tutte'a, kolorowania i dualność algebraiczna}
\label{ch:tg-tutte}

Graf planarny daje jednocześnie obraz geometryczny i zbiór danych
kombinatorycznych. Wielomian Tutte'a jest niezmiennikiem, który
łączy liczenie drzew rozpinających, kolorowania, przepływy
i dualność grafów płaskich. Zaczniemy od definicji przez
sumowanie po wszystkich podzbiorach krawędzi, a następnie
wyprowadzimy rekurencyjny algorytm obliczania tego wielomianu.
Bez dodatkowych uproszczeń liczba gałęzi rekurencji może
rosnąć wykładniczo względem liczby krawędzi.

\section{Usuwanie i kontrakcja krawędzi}
\label{sec:tg-deletion-contraction}

\begin{definicja}[Dwie operacje na krawędzi]
Niech $G=(V,E)$ będzie skończonym multigrafem, a $e\in E$.
W tym rozdziale $E$ jest zbiorem rozróżnialnych krawędzi:
równoległe krawędzie i pętle mają własne oznaczenia,
nawet gdy mają te same końce. Cykl może mieć długość
jeden (pętla) lub dwa (dwie równoległe krawędzie).
Przez $G\setminus e$ oznaczamy graf po usunięciu krawędzi $e$
bez usuwania jej końców. Jeśli $e$ nie jest pętlą,
przez $G/e$ oznaczamy graf powstały przez kontrakcję $e$:
utożsamiamy jej dwa końce i usuwamy $e$; pozostałe krawędzie
zachowują swoje wystąpienia, także jeśli po utożsamieniu
zamieniają się w pętle. Przy liczeniu wierzchołków grafu $G/e$
liczba wierzchołków maleje dokładnie o jeden.
\end{definicja}

\begin{definicja}[Rząd i wymiar przestrzeni cykli zbioru krawędzi]
Niech $G=(V,E)$ będzie skończonym multigrafem, a $A\subseteq E$.
Przez $k_G(A)$ oznaczamy liczbę składowych grafu $(V,A)$,
łącznie z izolowanymi wierzchołkami. Definiujemy
$r_G(A):=|V|-k_G(A)$ oraz
$\nu_G(A):=|A|-r_G(A)=|A|-|V|+k_G(A)$.
Liczba $\nu_G(A)$ jest wymiarem przestrzeni cykli grafu $(V,A)$
nad $\mathbb F_2$. Istotnie, las rozpinający $(V,A)$
otrzymujemy, usuwając pętle, zastępując każdą rodzinę
krawędzi równoległych jedną krawędzią i wybierając las
rozpinający uzyskanego grafu prostego na mocy twierdzenia
\ref{thm:tg-las-rozpinajacy}. Każdą wybraną krawędź
reprezentujemy przez jedną z jej pierwotnych kopii.
W każdej składowej tak otrzymanego lasu wybieramy korzeń $r$. Kolumny macierzy
incydencji nad $\mathbb F_2$ mają sumę współczynników
zero w każdej składowej (dla pętli kolumna jest zerowa).
Dla każdego wierzchołka $v$ różnego od korzenia suma
kolumn wzdłuż ścieżki od $r$ do $v$ daje $[v]+[r]$.
Wektory te tworzą bazę podprzestrzeni o sumie
współczynników zero w każdej składowej: współczynniki
poza korzeniami można wybrać dowolnie, a współczynniki
przy korzeniach są przez nie wyznaczone.
Stąd rząd macierzy wynosi $|V|-k_G(A)$, a wymiar
jej jądra to $\nu_G(A)$. Jądro składa się dokładnie z wektorów
charakterystycznych zbiorów krawędzi o parzystych
stopniach (w pętli liczymy dwie incydencje).
Aby utożsamić je z przestrzenią rozpiętą przez cykle,
ustalmy opisany wyżej las rozpinający $F$.
Każda krawędź $e\in A\setminus E(F)$ wyznacza cykl
$C_e$ złożony z $e$ i jedynej drogi w lesie między
jej końcami (dla pętli droga jest pusta).
Wektor $\chi_{E(C_e)}$ należy do jądra i ma
współczynnik jeden przy $e$ oraz zero przy wszystkich
pozostałych krawędziach spoza lasu.
Dla wektora $z$ z jądra odejmujemy sumę
$\sum_{e\in A\setminus E(F)}z_e\chi_{E(C_e)}$.
Pozostaje element jądra wsparty na lesie,
który musi być zerowy: niepusty nośnik
w lesie miałby wierzchołek stopnia jeden.
Te same współrzędne poza lasem dowodzą niezależności
wskazanych wektorów. Zatem tworzą one bazę
jądra także dla multigrafów. Argument jest odpowiednikiem
dowodu twierdzenia \ref{thm:tg-rank} dla multigrafów
nad $\mathbb F_2$; obejmuje także $V=\varnothing$,
gdy obie przestrzenie są zerowe.
W szczególności $\nu_G(A)=0$ dokładnie wtedy, gdy
$(V,A)$ jest lasem. Istotnie, las rozpinający tego
podgrafu ma $|V|-k_G(A)$ krawędzi, zatem $\nu_G(A)$
liczy krawędzie spoza tego lasu. Każda dodatkowa
krawędź zamyka cykl wraz z jedyną drogą w drzewie
między jej końcami (dla pętli droga jest pusta,
a dla krawędzi równoległych cykl może mieć długość dwa).
\end{definicja}

\begin{definicja}[Wielomian Tutte'a]
\label{def:tg-tutte}
Dla dowolnego skończonego multigrafu $G=(V,E)$ definiujemy
\[
T_G(x,y):=\sum_{A\subseteq E}
 (x-1)^{k_G(A)-k_G(E)}
 (y-1)^{|A|-|V|+k_G(A)}\in\mathbb Z[x,y].
\]
Wykładniki są nieujemne: dodawanie krawędzi nie zwiększa liczby
składowych, a $|A|-|V|+k_G(A)$ jest wymiarem przestrzeni cykli.
Dla grafu bez krawędzi mamy $T_G=1$, niezależnie od liczby
wierzchołków izolowanych.
\end{definicja}

\begin{twierdzenie}[Rekurencja usuwania i kontrakcji]
\label{thm:tg-tutte-recursion}
Niech $G$ będzie skończonym multigrafem, a $e\in E(G)$.
Wtedy zachodzą następujące tożsamości:
\begin{enumerate}[label=(\roman*)]
\item jeśli $e$ nie jest ani mostem, ani pętlą, to
$T_G=T_{G\setminus e}+T_{G/e}$;
\item jeśli $e$ jest mostem, to $T_G=x\,T_{G/e}$;
\item jeśli $e$ jest pętlą, to $T_G=y\,T_{G\setminus e}$.
\end{enumerate}
Wraz z wartością $T=1$ dla grafu bez krawędzi reguły te
jednoznacznie określają wielomian Tutte'a.
\end{twierdzenie}

\begin{proof}
Podzielmy sumę definiującą $T_G$ na dwa rodzaje składników:
podzbiorów $A\subseteq E\setminus\{e\}$ i podzbiorów
postaci $A=B\cup\{e\}$, gdzie $B\subseteq E\setminus\{e\}$.

Gdy $e$ nie jest mostem ani pętlą, liczba składowych całego
grafu nie zmienia się przy usunięciu ani przy kontrakcji $e$.
Składniki nie zawierające $e$ są zatem dokładnie składnikami
$T_{G\setminus e}$. Dla składników zawierających $e$ mamy
$k_G(B\cup\{e\})=k_{G/e}(B)$: utożsamiamy końce
krawędzi należącej do jednej składowej $(V,B\cup\{e\})$,
co zachowuje wszystkie składowe i nie łączy różnych. Zarazem
$|V(G/e)|=|V(G)|-1$, więc
\[
|B\cup\{e\}|-|V(G)|+k_G(B\cup\{e\})
=|B|-|V(G/e)|+k_{G/e}(B).
\]
Składniki z $e$ dają zatem dokładnie $T_{G/e}$.

Jeśli $e$ jest mostem, dla każdego $B\subseteq E\setminus\{e\}$
zachodzi $k_G(B)=k_G(B\cup\{e\})+1$:
końce mostu leżą w różnych składowych $(V,B)$,
gdyż droga między nimi w $B$ przeczyłaby temu,
że usunięcie $e$ rozspaja jego składową w $G$.
Oznaczmy $k=k_G(B)$, $c=k_G(E)$ i $n=|V|$.
Po kontrakcji $e$ mamy $k_{G/e}(B)=k-1$,
$k_{G/e}(E\setminus\{e\})=c$ oraz $|V(G/e)|=n-1$.
Dwa składniki związane z $B$ mają ten sam wykładnik
$|B|-n+k$ przy $y-1$, natomiast wykładniki przy $x-1$
wynoszą $k-c$ i $k-c-1$. Ich suma jest równa
\[
 x(x-1)^{k-c-1}(y-1)^{|B|-n+k}.
\]
Po wyłączeniu $x$ jest to dokładnie składnik $T_{G/e}$
odpowiadający $B$. Sumowanie po $B$ daje $T_G=xT_{G/e}$.
Jeżeli $e$ jest pętlą, to dla każdego
$B\subseteq E\setminus\{e\}$ zachodzą równości
$k_G(B\cup\{e\})=k_G(B)$ oraz
$\nu_G(B\cup\{e\})=\nu_G(B)+1$.
Ponadto $k_G(E)=k_{G\setminus e}(E\setminus\{e\})$.
Dwa odpowiadające sobie składniki mają więc sumę
\[
(x-1)^{k_G(B)-k_G(E)}(y-1)^{\nu_G(B)}
\bigl(1+(y-1)\bigr).
\]
Czynnik w nawiasie jest równy $y$, a pozostały
iloczyn jest składnikiem $T_{G\setminus e}$
odpowiadającym $B$. Sumowanie daje
$T_G=yT_{G\setminus e}$.

Każda z powyższych rekurencji zmniejsza liczbę krawędzi
o jeden, a więc pozwala obliczyć wielomian po skończonej
liczbie kroków. Jednoznaczność wynika z indukcji:
jeśli dwa przyporządkowania grafom wielomianów spełniają
te reguły i zgadzają się dla grafów bez krawędzi,
to dla grafu o $m>0$ krawędziach wybieramy tę samą
krawędź $e$ w obu przyporządkowaniach.
Na mocy założenia indukcyjnego wartości dla grafów
o $m-1$ krawędziach są równe, więc odpowiednia
reguła daje równe wartości także dla $G$.
Ponieważ suma z definicji \ref{def:tg-tutte} spełnia
wszystkie trzy reguły dla każdej dopuszczalnej krawędzi,
wynik obliczenia nie zależy od kolejności ich wyboru.
\end{proof}

\begin{wniosek}[Nieujemność współczynników i rozłączna suma grafów]
\label{cor:tg-tutte-product}
Współczynniki $T_G$ w bazie jednomianów $x^iy^j$
są nieujemnymi liczbami całkowitymi.
Jeśli $G=G_1\sqcup G_2$ jest rozłączną sumą grafów,
to $T_G=T_{G_1}T_{G_2}$.
Dodanie izolowanych wierzchołków nie zmienia $T_G$.
\end{wniosek}
\begin{proof}
Nieujemność wynika przez indukcję względem liczby
krawędzi z twierdzenia \ref{thm:tg-tutte-recursion}:
w przypadku bazowym wielomian jest równy $1$,
a w każdym kroku dodajemy dwa wielomiany o nieujemnych
współczynnikach albo mnożymy taki wielomian przez $x$
lub $y$.

Każdy podzbiór $A\subseteq E(G)$ rozkłada się
jednoznacznie na rozłączną sumę $A_1\sqcup A_2$,
gdzie $A_i\subseteq E(G_i)$. Mamy
\[
k_G(A)=k_{G_1}(A_1)+k_{G_2}(A_2),\qquad
\nu_G(A)=\nu_{G_1}(A_1)+\nu_{G_2}(A_2),
\]
ponieważ liczby składowych, wierzchołków i krawędzi
sumują się dla rozłącznych grafów.
Tak samo $k_G(E)=k_{G_1}(E(G_1))+k_{G_2}(E(G_2))$.
Składnik definicji \ref{def:tg-tutte} jest zatem
iloczynem odpowiednich składników dla $G_1,G_2$.
Sumowanie niezależnie po $A_1$ i $A_2$ daje
wzór na iloczyn. Graf bez krawędzi ma wielomian $1$,
co dowodzi ostatniej tezy.
\end{proof}

\begin{przyklad}[Drzewo, cykl i trzy krawędzie równoległe]
Jeżeli $G$ jest drzewem o $n\ge1$ wierzchołkach,
każda krawędź jest mostem. Kontraktując je kolejno,
otrzymujemy $T_G(x,y)=x^{n-1}$.

Dla trójkąta wybierzmy krawędź $e$.
Graf $C_3\setminus e$ jest ścieżką o dwóch krawędziach,
z wielomianem $x^2$. Graf $C_3/e$ składa się z dwóch
równoległych krawędzi; ich wielomian to $x+y$.
Stąd
$T_{C_3}=x^2+x+y$.
Dla $n=2$ przyjmujemy, że $C_2$ jest multigrafem
z dwiema równoległymi krawędziami. Usunięcie jednej
pozostawia most, a kontrakcja pozostawia jedną pętlę,
więc $T_{C_2}=x+y$. Dla $n\ge3$ usunięcie krawędzi
cyklu daje drzewo z $n-1$ krawędziami, a kontrakcja
daje $C_{n-1}$. Stąd $T_{C_n}=x^{n-1}+T_{C_{n-1}}$
i przez indukcję
\[
 T_{C_n}(x,y)=\sum_{j=1}^{n-1}x^j+y,\qquad n\ge2.
\]
Dla grafu theta złożonego z dwóch wierzchołków
i trzech równoległych krawędzi usunięcie jednej
krawędzi daje $C_2$, a kontrakcja daje dwie pętle
przy jednym wierzchołku, o wielomianie $y^2$.
Zatem $T_{\Theta}=T_{C_2}+y^2=x+y+y^2$.
\end{przyklad}

\begin{center}
\begin{tikzpicture}[scale=.90,>=Stealth,line cap=round]
 \begin{scope}[shift={(-3.3,0)}]
  \node[tgBlueVertex] (a) at (-1,-.7) {$a$};
  \node[tgBlueVertex] (b) at (1,-.7) {$b$};
  \node[tgBlueVertex] (c) at (0,1.05) {$c$};
  \draw[tgEdge] (a)--(b)--(c);
  \draw[tgAccent] (c)-- node[midway,left,font=\scriptsize] {$e$} (a);
  \node[text=tgNavy,font=\small] at (0,-1.45) {$T_{C_3}$};
 \end{scope}
 \node[text=tgNavy,font=\large] at (-1.6,0) {$=$};
 \begin{scope}[shift={(.1,0)}]
  \node[tgBlueVertex] (a) at (-.7,-.7) {};
  \node[tgBlueVertex] (b) at (.7,-.7) {};
  \node[tgBlueVertex] (c) at (0,1.05) {};
  \draw[tgTree] (a)--(b)--(c);
  \node[text=tgNavy,font=\footnotesize] at (0,1.65) {$C_3\setminus e$};
  \node[text=tgNavy,font=\small] at (0,-1.45) {$x^2$};
 \end{scope}
 \node[text=tgNavy,font=\large] at (2.0,0) {$+$};
 \begin{scope}[shift={(4.0,0)}]
  \node[tgBlueVertex] (a) at (-.7,0) {};
  \node[tgBlueVertex] (b) at (.7,0) {};
  \draw[tgBlue,line width=2pt] (a) to[bend left=38] (b);
  \draw[tgBlue,line width=2pt] (a) to[bend right=38] (b);
  \node[text=tgNavy,font=\footnotesize] at (0,1.65) {$C_3/e$};
  \node[text=tgNavy,font=\small] at (0,-1.45) {$x+y$};
 \end{scope}
\end{tikzpicture}

{\footnotesize Rekurencja Tutte'a dla trójkąta: usunięcie krawędzi
daje drzewo, kontrakcja daje dwie krawędzie równoległe.}
\end{center}

\section{Drzewa rozpinające jako specjalizacja wielomianu}
\label{sec:tg-tutte-trees}

\begin{twierdzenie}[Wartość $T_G(1,1)$]
Niech $G$ będzie skończonym spójnym multigrafem.
Wtedy $T_G(1,1)=\tau(G)$, gdzie $\tau(G)$ oznacza
liczbę drzew rozpinających, z uwzględnieniem krotności
krawędzi równoległych.
\end{twierdzenie}

\begin{proof}
Dla grafu spójnego $k_G(E)=1$. W sumie definiującej
$T_G(1,1)$ składnik odpowiadający $A\subseteq E$
nie znika wtedy i tylko wtedy, gdy obie potęgi zerowej
podstawy mają wykładnik zero, czyli
$k_G(A)=1$ oraz $|A|-|V|+k_G(A)=0$.
Te warunki oznaczają, że $(V,A)$ jest spójnym grafem
bez cykli, czyli drzewem rozpinającym. Każdy taki
składnik jest równy jeden ($0^0$ rozumiemy tu
jako wartość jednomianu z wykładnikiem zero).
Suma liczy dokładnie drzewa rozpinające.
\end{proof}

\begin{uwaga}
Twierdzenie \ref{thm:tg-kirchhoff} daje tę samą liczbę
z wyznacznika minora laplasjanu. Aby zastosować jego
wersję dla grafów prostych do spójnego multigrafu,
usuwamy pętle (nie należą do żadnego drzewa),
a każdą rodzinę $s$ równoległych krawędzi zastępujemy
jedną krawędzią o wadze $s$. Każdemu drzewu $T$ grafu prostego odpowiada dokładnie
$\prod_{e\in E(T)}w_e$ różnych drzew pierwotnego
multigrafu: dla każdej krawędzi drzewa niezależnie
wybieramy jedną z jej $w_e$ kopii.
Laplasjan $BB^T$ multigrafu ma zerowy wkład od pętli,
a $s$ równoległych krawędzi wnosi ten sam składnik
co jedna krawędź o wadze $s$. Zatem ważona wersja
twierdzenia Kirchhoffa daje żądaną liczbę również
dla multigrafu, z konwencją pustego minora dla $|V|=1$.
Mamy zatem dwie niezależne metody
obliczania $\tau(G)$: algebraiczną (wyznacznik) i
kombinatoryczną (wielomian Tutte'a). Dla grafu niespójnego
warunki przetrwania składnika przy $(1,1)$ to
$k_G(A)=k_G(E)$ i $\nu_G(A)=0$. Pierwszy oznacza,
że podgraf $(V,A)$ jest spójny na każdej składowej
$G$, ponieważ jego składowe nie mogą łączyć różnych
składowych $G$. Drugi oznacza brak cykli.
Zatem $T_G(1,1)$ liczy maksymalne lasy rozpinające,
po jednym drzewie rozpinającym na każdą składową.
Dla grafu zerowego jedynym takim lasem jest graf pusty.
\end{uwaga}

\begin{stwierdzenie}[Trzy dalsze specjalizacje]
\label{prop:tg-tutte-specializations}
Dla skończonego multigrafu $G=(V,E)$ zachodzą równości:
\[
T_G(2,2)=2^{|E|},\qquad
T_G(2,1)=N_{\mathrm{las}}(G),\qquad
T_G(1,2)=N_{\mathrm{rozp}}(G),
\]
gdzie $N_{\mathrm{las}}(G)$ oznacza liczbę podzbiorów
krawędzi niezawierających cyklu (w tym podzbiór pusty),
a $N_{\mathrm{rozp}}(G)$ liczbę podzbiorów $A\subseteq E$,
dla których podgraf $(V,A)$ jest spójny na każdej
składowej $G$.
\end{stwierdzenie}

\begin{proof}
Dla $x=y=2$ w definicji każdy składnik ma wartość jeden,
a podzbiorów $E$ jest $2^{|E|}$. Przy $x=2$, $y=1$
pozostają dokładnie składniki z $\nu_G(A)=0$, czyli
podgrafy $(V,A)$ będące lasami. Każdy taki składnik
wnosi jeden. Pętle nigdy nie należą do lasu i słusznie
nie pojawiają się w tej sumie.
Przy $x=1$, $y=2$ pozostają dokładnie składniki
z $k_G(A)=k_G(E)$. Podgraf $(V,A)$ nie może łączyć
różnych składowych $G$, zatem równość liczb składowych
zachodzi wtedy i tylko wtedy, gdy żadna składowa $G$
nie została rozspójniona przez usunięcie krawędzi
spoza $A$. Każdy taki składnik wnosi jeden,
co daje $N_{\mathrm{rozp}}(G)$.
W szczególności dla spójnego $G$ wartość $T_G(1,2)$
liczy spójne podgrafy rozpinające, nie tylko drzewa.
\end{proof}

\section{Wielomian chromatyczny i jego wzór przez Tutte'a}
\label{sec:tg-chromatic}

\begin{uwaga}[Konwencje dla grafu zerowego i pętli]
W tej części dopuszczamy również graf zerowy: ma on dokładnie
jedno kolorowanie pustą funkcją i $P_{G_0}(q)=1$, ale jego
liczbę chromatyczną definiujemy osobno jako zero. Jeśli graf
posiada pętlę, nie istnieje poprawne kolorowanie żadną liczbą
kolorów, ponieważ wierzchołek nie może mieć koloru różnego
od samego siebie. Zatem jego wielomian chromatyczny jest
wielomianem zerowym. Krawędzie równoległe nie zmieniają
liczby poprawnych kolorowań.
\end{uwaga}

\begin{definicja}[Wielomian chromatyczny]
Niech $G=(V,E)$ będzie skończonym multigrafem.
Dla $q\in\mathbb N$, $q\ge1$, przez $P_G(q)$ oznaczamy
liczbę poprawnych kolorowań wierzchołków kolorami
$1,\ldots,q$, tzn. funkcji $c:V\to\{1,\ldots,q\}$
spełniających $c(u)\ne c(v)$ dla każdej krawędzi $uv$.
Gdy $G$ ma pętlę, przyjmujemy $P_G(q)=0$.
\end{definicja}

\begin{twierdzenie}[Wzór włączeń i wyłączeń]
\label{thm:tg-chromatic-inclusion}
Dla dowolnego skończonego multigrafu $G=(V,E)$ i każdej
liczby całkowitej $q\ge1$ zachodzi
\[
 P_G(q)=\sum_{A\subseteq E}(-1)^{|A|}q^{k_G(A)}.
\]
Prawa strona jest wielomianem w $q$ o współczynnikach
całkowitych, wyznaczającym liczbę kolorowań dla każdego $q\ge1$.
Jest to jedyny taki wielomian, ponieważ różnica dwóch
wielomianów zgodnych dla wszystkich dodatnich liczb
całkowitych ma nieskończenie wiele pierwiastków,
a więc jest wielomianem zerowym.
\end{twierdzenie}

\begin{proof}
Wszystkich funkcji kolorujących wierzchołki jest $q^{|V|}$.
Dla każdej krawędzi $e=uv$ oznaczmy przez $B_e$
zbiór funkcji, które mają $c(u)=c(v)$.
Poprawne kolorowania to dokładnie funkcje nienależące
do żadnego z $B_e$. Z zasady włączeń i wyłączeń ich liczba
to $\sum_{A\subseteq E}(-1)^{|A|}\bigl|\bigcap_{e\in A}B_e\bigr|$.
Dla ustalonego $A$ warunki $c(u)=c(v)$ na krawędziach $A$
wymuszają stałość kolorowania na każdej składowej $(V,A)$.
Składowych jest $k_G(A)$ i każda może niezależnie otrzymać
jeden z $q$ kolorów. Dlatego
$\bigl|\bigcap_{e\in A}B_e\bigr|=q^{k_G(A)}$.
Podstawienie kończy dowód, również w przypadku pętli.
\end{proof}

\begin{twierdzenie}[Wielomian chromatyczny jako specjalizacja Tutte'a]
\label{thm:tg-chrom-tutte}
Niech $G=(V,E)$ będzie skończonym multigrafem o $n=|V|$
wierzchołkach i $c=k_G(E)$ składowych. Wtedy
\[
P_G(q)=(-1)^{n-c}q^c\,T_G(1-q,0)
\]
jako tożsamość wielomianów w $\mathbb Z[q]$.
\end{twierdzenie}

\begin{proof}
Podstawiając $x=1-q$, $y=0$ w definicji Tutte'a,
otrzymujemy
\[
(-1)^{n-c}q^c T_G(1-q,0)
=\sum_{A\subseteq E}
 (-1)^{n-c+k_G(A)-c+|A|-n+k_G(A)}q^{k_G(A)}.
\]
Wykładnik znaku $-1$ różni się od $|A|$ o liczbę parzystą
$2(k_G(A)-c)$, dlatego każdy składnik ma postać
$(-1)^{|A|}q^{k_G(A)}$. Poprzednie twierdzenie
identyfikuje sumę z $P_G(q)$.
\end{proof}

\begin{wniosek}[Rekurencja wielomianu chromatycznego]
\label{cor:tg-chrom-recursion}
Jeśli $e$ nie jest pętlą, to
\[
P_G(q)=P_{G\setminus e}(q)-P_{G/e}(q).
\]
Dla grafu bez krawędzi $P_G(q)=q^{|V|}$,
a dla grafu z pętlą $P_G(q)=0$.
\end{wniosek}
\begin{proof}
Dla każdej dodatniej liczby całkowitej $q$ dzielimy
poprawne kolorowania $G\setminus e$ na dwie rozłączne
rodziny według tego, czy końce $e$ mają różne kolory,
czy ten sam kolor. Pierwszą rodzinę stanowią dokładnie
poprawne kolorowania $G$. Drugą utożsamiamy bijektywnie
z poprawnymi kolorowaniami $G/e$: końce $e$ otrzymują
wspólny kolor i zostają utożsamione, a wszystkie
pozostałe ograniczenia kolorowania są zachowane.
Jeżeli po kontrakcji powstaje pętla, druga rodzina
jest pusta, zgodnie z konwencją $P_{G/e}=0$.
Zliczanie obu rodzin daje rekurencję dla każdego
całkowitego $q\ge1$, a więc tożsamość wielomianów
na mocy twierdzenia \ref{thm:tg-chromatic-inclusion}.
Dwie wartości początkowe wynikają z definicji.
\end{proof}

\begin{przyklad}[Kolorowania trójkąta]
Ponieważ $T_{C_3}(x,y)=x^2+x+y$, mamy
\[
P_{C_3}(q)=q\bigl((1-q)^2+(1-q)\bigr)
=q(q-1)(q-2).
\]
Dla $q\ge2$ odpowiada to bezpośredniemu rachunkowi:
pierwszy wierzchołek kolorujemy na $q$ sposobów,
drugi na $q-1$, a trzeci na $q-2$ sposoby.
Dla $q=1$ nie istnieje poprawne kolorowanie
i otrzymany wielomian również ma wartość zero.
W szczególności $\chi(C_3)=3$.
\end{przyklad}

\section{Dualność wielomianu Tutte'a}
\label{sec:tg-tutte-duality}

\begin{twierdzenie}[Dualność Tutte'a]
\label{thm:tg-tutte-dual}
Niech $G\subset S^2$ będzie skończonym spójnym grafem
płaskim, a $G^*$ jego dualnym multigrafem. Wtedy
\[
 T_{G^*}(x,y)=T_G(y,x).
\]
\end{twierdzenie}

\begin{proof}
Niech $n=|V(G)|$, $m=|E(G)|$, $f=|V(G^*)|$.
Spójne zanurzenie w sferze jest komórkowe:
w dowodzie twierdzenia \ref{thm:tg-euler-plane}
jego regularne otoczenie jest spójne, a każda składowa
dopełnienia tego otoczenia jest dyskiem. Dołączenie
półotwartego kołnierza daje otwarty dysk, czyli ścianę.
Możemy zatem stosować konstrukcję grafu dualnego
z rozdziału \ref{ch:tg-powierzchnie}.
Dla $A\subseteq E(G)$ oznaczmy
$A^\perp=\{e^*:e\in E(G)\setminus A\}$.
Symbol $A^\perp$ oznacza tu dopełnienie zbioru
krawędzi przeniesione na graf dualny, a nie
dopełnienie ortogonalne podprzestrzeni.
Zauważmy, że składowe podgrafu $(V(G^*),A^\perp)$
odpowiadają ścianom grafu $(V(G),A)$ zanurzonego w sferze.
Istotnie, usuwamy kolejno krawędzie spoza $A$, nie usuwając
wierzchołków. Usunięcie wnętrza krawędzi łączy dokładnie
dwie przyległe ściany, jeżeli są różne; jeśli już należą
do tej samej ściany, nie zmienia liczby ścian.
Jest to lokalny argument z dowodu wniosku
\ref{cor:tg-bridge-face}, który działa również dla grafu
niespójnego i dla pętli. W grafie dualnym w tym samym
kroku dokładamy krawędź $e^*$, łącząc dokładnie składowe
odpowiadające tym ścianom. Na początku każda ściana $G$
odpowiada pojedynczemu wierzchołkowi dualnemu; indukcja
po usuwanych krawędziach daje żądaną odpowiedniość.

Z twierdzenia \ref{thm:tg-euler-plane}, zastosowanego
do $(V(G),A)$ (z uwzględnieniem izolowanych wierzchołków),
wynika
\[
 k_{G^*}(A^\perp)=1+k_G(A)-n+|A|=1+\nu_G(A).
\]
Ponieważ $G$ i $G^*$ są spójne, a $f=m-n+2$,
otrzymujemy również
\[
\begin{aligned}
 \nu_{G^*}(A^\perp)
 &=m-|A|-f+k_{G^*}(A^\perp)\\
 &=m-|A|-(m-n+2)+1+k_G(A)-n+|A|\\
 &=k_G(A)-1.
\end{aligned}
\]
Zatem składnik sumy definiującej $T_{G^*}(x,y)$
odpowiadający $A^\perp$ jest równy
$(x-1)^{\nu_G(A)}(y-1)^{k_G(A)-1}$.
Przyporządkowanie $A\mapsto A^\perp$ jest bijekcją
między zbiorami wszystkich podzbiorów krawędzi obu grafów.
Sumowanie daje $T_{G^*}(x,y)=T_G(y,x)$.
Argument obejmuje także $m=0$: obydwa grafy mają po jednym
wierzchołku, a jedyny składnik sumy jest równy jeden.
\end{proof}

\begin{przyklad}[Trójkąt i trzy krawędzie równoległe]
Na sferze grafem dualnym trójkąta $C_3$ jest multigraf złożony
z dwóch wierzchołków połączonych trzema równoległymi
krawędziami. Zatem
$T_{C_3}=x^2+x+y$ oraz
$T_{C_3^*}=x+y+y^2=T_{C_3}(y,x)$.
Dualność zamienia wykładniki opisujące wymiar przestrzeni
cykli i różnicę liczby składowych.
\end{przyklad}

\section{Kolorowanie grafu dualnego a przepływy}
\label{sec:tg-flows-duality}

\begin{definicja}[Przepływ nigdzie niezerowy modulo $q$]
Niech $G=(V,E)$ będzie skończonym multigrafem z wybraną
orientacją krawędzi, a $q\ge2$ liczbą całkowitą.
Przez $B$ oznaczamy macierz incydencji wierzchołek--krawędź
z rozdziału \ref{ch:tg-algebra}: kolumna krawędzi
od $u$ do $v$ jest równa $[v]-[u]$, więc dla pętli
jest zerowa, a krawędzie równoległe mają osobne kolumny.
Współczynniki macierzy redukujemy modulo $q$.
\emph{Przepływem o wartościach w $\mathbb Z/q\mathbb Z$}
nazywamy wektor $j\in(\mathbb Z/q\mathbb Z)^E$
spełniający $Bj=0$. Jest on \emph{nigdzie niezerowy}, jeżeli
$j_e\ne0$ dla każdej krawędzi. Liczbę takich przepływów
oznaczamy przez $F_G(q)$.
Liczba ta nie zależy od pomocniczej orientacji: odwrócenie
krawędzi zmienia znak jej kolumny, a jednoczesna zamiana
$j_e$ na $-j_e$ daje bijekcję przepływów, zachowującą
warunek niezerowości.
\end{definicja}

Warunek $Bj=0$ jest warunkiem zachowania na wszystkich
wierzchołkach, w odróżnieniu od przepływu między źródłem
i ujściem z rozdziału \ref{ch:tg-menger}.
Tutaj współczynniki należą do pierścienia
$\mathbb Z/q\mathbb Z$ i nie ma ograniczeń
przepustowości ani maksymalizowanej wartości.
Dla $q=2$ warunek ten opisuje przestrzeń cykli
z rozdziału \ref{ch:tg-algebra}.

\begin{twierdzenie}[Wielomian przepływów jako specjalizacja Tutte'a]
\label{thm:tg-flow-tutte}
Niech $G=(V,E)$ będzie dowolnym skończonym multigrafem.
Dla każdej liczby całkowitej $q\ge2$ zachodzi
\[
F_G(q)=\sum_{A\subseteq E}(-1)^{|E|-|A|}q^{\nu_G(A)}
=(-1)^{\nu_G(E)}T_G(0,1-q).
\]
W szczególności liczba przepływów nigdzie niezerowych
jest wartością wielomianu o współczynnikach całkowitych.
Wielomian ten jest jednoznaczny, ponieważ dwa wielomiany
zgodne dla wszystkich liczb całkowitych $q\ge2$ mają
różnicę z nieskończenie wieloma pierwiastkami.
Wynik nie wymaga planarności ani spójności grafu.
\end{twierdzenie}

\begin{proof}
Niech $R=\mathbb Z/q\mathbb Z$.
Dla $A\subseteq E$ oznaczmy przez $B_A$
macierz incydencji grafu $(V,A)$ nad $R$.
Pokażemy najpierw, że
\[
|\ker B_A|=q^{\nu_G(A)}.
\]
Wybierzmy las rozpinający $T$ grafu $(V,A)$,
po jednym drzewie w każdej składowej.
Dla każdej krawędzi $e\in A\setminus E(T)$
tworzymy całkowity cykl fundamentalny $z^{(e)}$:
krawędź $e$ z jej wybraną orientacją oraz ścieżkę
w $T$ od końca $e$ do jej początku, ze znakami
zgodnymi z kierunkiem przejścia.
Dla pętli ścieżka jest pusta.
Brzeg tego łańcucha jest zerowy, a jego współczynnik
przy $e$ wynosi jeden i przy każdej innej krawędzi
spoza lasu wynosi zero.

Odwzorowanie
\[
\ker B_A\longrightarrow R^{A\setminus E(T)},\qquad
j\longmapsto (j_e)_{e\in A\setminus E(T)}
\]
jest izomorfizmem modułów. Jest surjektywne,
gdyż dowolna rodzina $(a_e)$ jest obrazem
$\sum_e a_e z^{(e)}$, z redukcją współczynników modulo $q$.
Jest injektywne: przepływ znikający poza lasem
ma przy krawędzi incydentnej z liściem współczynnik
zerowy na mocy równania zachowania w tym liściu.
Usuwając takie krawędzie kolejno, otrzymujemy
zerowy przepływ. Izolowane wierzchołki nie nakładają
dodatkowych warunków.
Ponieważ $|A\setminus E(T)|=|A|-|V|+k_G(A)=\nu_G(A)$,
otrzymujemy żądaną liczność. Nie korzystaliśmy
z dzielenia przez elementy pierścienia $R$, więc
dowód jest ważny również dla złożonych $q$.

W zbiorze wszystkich przepływów na $G$ niech $Z_e$
oznacza zbiór przepływów z zerowym współczynnikiem
przy krawędzi $e$. Dla $S\subseteq E$ przecięcie
$\bigcap_{e\in S}Z_e$ jest w bijekcji z $\ker B_{E\setminus S}$:
ograniczamy przepływ do pozostałych krawędzi,
a odwrotnie przedłużamy go przez zero.
Z zasady włączeń i wyłączeń wynika
\[
F_G(q)=\sum_{S\subseteq E}(-1)^{|S|}
q^{\nu_G(E\setminus S)}.
\]
Podstawienie $A=E\setminus S$ daje pierwszą równość.

Oznaczmy $n=|V|$, $m=|E|$ i $c=k_G(E)$.
W definicji \ref{def:tg-tutte} podstawiamy $x=0$,
$y=1-q$. Po pomnożeniu przez $(-1)^{m-n+c}$
składnik odpowiadający $A$ ma znak o wykładniku
\[
m-n+c+k_G(A)-c+\nu_G(A)
=m+|A|-2n+2k_G(A).
\]
Wykładnik ten ma tę samą parzystość co $m-|A|$.
Pozostałym czynnikiem jest $q^{\nu_G(A)}$.
Sumowanie daje drugą równość.
Dla $E=\varnothing$ jedyny przepływ jest pusty
i oba wzory dają jeden, także gdy $V=\varnothing$.
\end{proof}

\begin{wniosek}[Grupa współczynników a liczba przepływów]
Liczba przepływów nigdzie niezerowych o wartościach
w skończonej grupie abelowej $M$ zależy wyłącznie
od jej rzędu $q=|M|\ge2$ i wynosi $F_G(q)$.
\end{wniosek}
\begin{proof}
Równania zachowania zapisujemy addytywnie w $M$.
W poprzednim dowodzie izomorfizm przez współczynniki
poza lasem działa również dla grup abelowych:
całkowite współczynniki cyklu oznaczają odpowiednie
wielokrotności elementów $M$, a eliminacja liści
wymaga wyłącznie dodawania i odejmowania.
Otrzymujemy zatem dokładnie $|M|^{\nu_G(A)}$
przepływów na każdym podgrafie $(V,A)$.
Ta sama zasada włączeń i wyłączeń daje $F_G(q)$.
\end{proof}

\begin{wniosek}[Rekurencja wielomianu przepływów]
\label{cor:tg-flow-recursion}
Dla krawędzi $e$ skończonego multigrafu $G$ zachodzi
\[
F_G(q)=
\begin{cases}
0,& e\text{ jest mostem},\\
(q-1)F_{G\setminus e}(q),& e\text{ jest pętlą},\\
F_{G/e}(q)-F_{G\setminus e}(q),&
e\text{ nie jest mostem ani pętlą}.
\end{cases}
\]
Wartością początkową jest $F_G(q)=1$ dla grafu bez krawędzi.
\end{wniosek}
\begin{proof}
Dla mostu wzór wynika z twierdzeń
\ref{thm:tg-flow-tutte} i \ref{thm:tg-tutte-recursion},
ponieważ $T_G(0,y)=0$.
Można to również uzasadnić bezpośrednio: sumujemy
równania zachowania po wszystkich wierzchołkach jednej
z dwóch składowych powstałych przez usunięcie mostu
z jego pierwotnej składowej. Wkłady krawędzi
wewnętrznych znoszą się, a jedynym pozostałym
wkładem jest $\pm j_e$, więc $j_e=0$.

Dla pętli kolumna macierzy incydencji jest zerowa,
zatem każdy przepływ nigdzie niezerowy na
$G\setminus e$ można przedłużyć niezależnie
na $q-1$ sposobów przez wybór $j_e\ne0$.

Niech teraz $e$ nie będzie ani mostem, ani pętlą.
Oznaczmy $\nu=\nu_G(E(G))$.
Usunięcie $e$ nie zmienia liczby składowych,
a kontrakcja zmniejsza zarówno liczbę wierzchołków,
jak i krawędzi o jeden, nie zmieniając liczby składowych.
Dlatego
\[
\nu_{G\setminus e}(E(G\setminus e))=\nu-1,\qquad
\nu_{G/e}(E(G/e))=\nu.
\]
Mnożymy rekurencję Tutte'a dla $(x,y)=(0,1-q)$
przez $(-1)^\nu$ i stosujemy twierdzenie
\ref{thm:tg-flow-tutte} do obu mniejszych grafów.
Z uwagi na różnicę parzystości wykładników otrzymujemy
$F_G=F_{G/e}-F_{G\setminus e}$.
Wartość początkowa wynika z istnienia jedynego
pustego przepływu.
\end{proof}

\begin{uwaga}[Przepływy modulo dwa]
Nad $\mathbb Z/2\mathbb Z$ istnieje dokładnie jedna
możliwa rodzina niezerowych współczynników:
$j_e=1$ dla wszystkich $e\in E$.
Jej brzeg jest zerowy wtedy i tylko wtedy, gdy
każdy wierzchołek ma parzysty stopień, przy czym
pętla wnosi do stopnia dwa.
Zatem $F_G(2)$ jest równe jeden dla grafu o wszystkich
stopniach parzystych i zero w przeciwnym przypadku.
Nie zakładamy tutaj spójności grafu.
\end{uwaga}

\begin{twierdzenie}[Dualność kolorowań i przepływów]
\label{thm:tg-flow-color-dual}
Niech $G\subset S^2$ będzie skończonym spójnym grafem płaskim,
a $G^*$ jego grafem dualnym. Dla każdej liczby całkowitej $q\ge2$ zachodzi
\[
 P_{G^*}(q)=q\,F_G(q).
\]
\end{twierdzenie}

\begin{proof}
Ustalmy orientację sfery i pomocniczą orientację każdej
krawędzi $e$ grafu $G$. Niech $F_L(e)$ i $F_R(e)$ oznaczają
ściany po lewej i prawej stronie zorientowanej krawędzi.
Niech $R=\mathbb Z/q\mathbb Z$. Zbiór $q$ kolorów
utożsamiamy bijektywnie z $R$; warunek różnych kolorów
nie zależy od wyboru tej bijekcji. Zanurzenie jest
komórkowe, jak uzasadniono w dowodzie twierdzenia
\ref{thm:tg-tutte-dual}. Z dowolnej funkcji
$c:F(G)\to R$ tworzymy wektor
\[
 D(c)_e=c(F_R(e))-c(F_L(e)).
\]
Z definicji operatora brzegu dwukomórek w rozdziale
\ref{ch:tg-powierzchnie} współczynnik $\partial_2 c$
przy $e$ jest równy $c(F_L(e))-c(F_R(e))$.
Zatem $D(c)=-\partial_2 c$, gdzie wszystkie współczynniki
redukujemy modulo $q$. Ponieważ $\partial_1\partial_2=0$
nad $\mathbb Z$, ta sama tożsamość zachodzi nad $R$,
a więc $BD(c)=0$. Wartość $D(c)_e$ jest niezerowa
dokładnie wtedy, gdy końce $e^*$ mają różne kolory.

Wykażemy, że każdy wektor $j\in R^E$ spełniający $Bj=0$
ma postać $D(c)$. Wybierzmy drzewo rozpinające $T$ grafu $G$.
Na mocy twierdzenia \ref{thm:tg-dual-spanning} krawędzie
dualne do krawędzi spoza $T$ tworzą drzewo rozpinające
$T^*$ grafu $G^*$. Wybieramy korzeń $T^*$ i przypisujemy
mu kolor $0$. Idąc wzdłuż drzewa od korzenia do kolejnych
ścian, wyznaczamy kolor nowej ściany przez dodanie lub
odjęcie $j_e$ tak, aby $D(c)_e=j_e$ na każdej krawędzi
spoza $T$. Jedyność drogi w drzewie gwarantuje, że nie
powstaną sprzeczne wymagania.

Wektor $w=j-D(c)$ spełnia $Bw=0$ i znika na wszystkich
krawędziach spoza $T$. Jeżeli $T$ ma krawędź, wybieramy
liść drzewa; równanie $Bw=0$ w tym wierzchołku ma postać
$\pm w_e=0$, ponieważ wkłady wszystkich innych krawędzi
są zerowe. Stąd $w_e=0$. Usuwamy liść i jego krawędź
z drzewa i powtarzamy rozumowanie, aż nie pozostanie
żadna krawędź. Gdy $T$ od początku nie ma krawędzi,
$w=0$ bez tego kroku. Otrzymujemy $j=D(c)$.
W dowodzie nie dzieliliśmy przez żadną liczbę w $R$,
więc obejmuje on także złożone wartości $q$.

Jeżeli $D(c)=D(c')$, różnica $c-c'$ jest stała wzdłuż
każdej krawędzi grafu dualnego. Jego spójność oznacza,
że różnica jest jedną stałą na wszystkich ścianach.
Odwrotnie, dodanie dowolnej wspólnej stałej nie zmienia
$D(c)$. Każdy przepływ ma więc dokładnie $q$ przeciwobrazów,
a dla przepływu nigdzie niezerowego wszystkie te funkcje
są poprawnymi kolorowaniami. Dla grafu bez krawędzi
jedyny przepływ jest pustym wektorem, a graf dualny ma
jeden wierzchołek i $q$ kolorowań, zgodnie ze wzorem.
Stąd $P_{G^*}(q)=qF_G(q)$.
\end{proof}

Ten sam wzór otrzymujemy algebraicznie ze specjalizacji
Tutte'a. Istotnie, dla spójnego grafu płaskiego o $n$
wierzchołkach, $m$ krawędziach i $f$ ścianach mamy
$\nu_G(E)=m-n+1=f-1$ na mocy wzoru Eulera.
Z twierdzeń \ref{thm:tg-chrom-tutte},
\ref{thm:tg-tutte-dual} i \ref{thm:tg-flow-tutte} wynika
\[
\begin{aligned}
P_{G^*}(q)
&=(-1)^{f-1}qT_{G^*}(1-q,0)\\
&=q(-1)^{\nu_G(E)}T_G(0,1-q)
=qF_G(q).
\end{aligned}
\]
Dowód konstrukcyjny daje dodatkowo informację,
które kolorowania odpowiadają danemu przepływowi:
tworzą one jedną klasę względem dodawania wspólnej stałej.

\begin{uwaga}[Ograniczenie dualności przepływów do sfery]
Dla spójnego grafu komórkowo zanurzonego
w zamkniętej orientowalnej powierzchni $\Sigma_g$
nadal mamy $D=-\partial_2$ oraz $\partial_1D=0$
nad $R=\mathbb Z/q\mathbb Z$.
Jednak na mocy uwagi \ref{rem:tg-surface-mod-q}
iloraz wszystkich przepływów (również tych
z zerowymi współrzędnymi) przez przepływy
otrzymane z funkcji na ścianach jest izomorficzny
z $R^{2g}$. Dla $g>0$ nie każdy przepływ
jest więc różnicą kolorów ścian.
Na przykład dla standardowego komórkowego
zanurzenia bukietu dwóch okręgów w torusie
istnieje jedna ściana, więc $D(c)=0$ dla
każdej funkcji $c$. Obie krawędzie są pętlami,
a więc operator incydencji jest zerowy i
dowolna para niezerowych elementów $R$
jest przepływem nigdzie niezerowym.
Graf dualny ma pętle i nie ma poprawnego
kolorowania. Równość z twierdzenia
\ref{thm:tg-flow-color-dual} nie przenosi się
zatem bez dodatkowych warunków na wyższy rodzaj.
\end{uwaga}

% Zachowujemy końcowy przykład i literaturę na tej samej stronie,
% zamiast pozostawiać osobną stronę z krótką bibliografią.
\pagebreak[4]
\begin{przyklad}[Przepływy na grafie theta]
Graf theta z trzema równoległymi krawędziami, zanurzony
komórkowo w sferze, ma jako graf dualny trójkąt. Dlatego
\[
 F_{\Theta}(q)=\frac{P_{C_3}(q)}{q}=(q-1)(q-2).
\]
Można to policzyć również bezpośrednio. Nadajmy wszystkim
trzem krawędziom orientację od $u$ do $v$.
Warunek przepływu to $j_1+j_2+j_3=0$ modulo $q$.
Wybieramy $j_1\ne0$ na $q-1$ sposobów, następnie
$j_2\ne0,-j_1$ na $q-2$ sposoby; wtedy
$j_3=-(j_1+j_2)$ jest jednoznacznie wyznaczone i niezerowe.
Zgodność z rachunkiem dualnym jest całkowita.
\end{przyklad}

\section*{Dalsza geometria i literatura}
Długości krawędzi prowadzą do metryki i laplasjanu, a zanurzenie
komórkowe do $\partial_2$ i układów rotacji. Wykorzystamy te dane
w triangulacjach, grafach Reeba, teorii Morse'a i grafach plumbingowych.

Literatura uzupełniająca: R. Diestel, \emph{Graph Theory}
(twierdzenia Kuratowskiego i Wagnera); B. Mohar i C. Thomassen,
\emph{Graphs on Surfaces} (zanurzenia, rotacje i grafy wstęgowe);
\emph{Handbook of the Tutte Polynomial and Related Topics}
(kolorowania i przepływy).

% ================================================================
% ITERACJA 5. CZESC TOPOLOGICZNA
% ================================================================

